\subsection{Módulo de Visualización y Gemelo Digital (VIS)}
\noindent\textbf{Ficha Técnica: VIS-008 - Visualización de Capa de Permisos de Trabajo}
\footnotesize\setlength{\tabcolsep}{4pt}\renewcommand{\arraystretch}{1.2}
\rowcolors{2}{tableBlue}{white}\begin{longtable}{|p{0.21\textwidth}|p{0.74\textwidth}|}\hline
\textbf{Número} & VIS-008 \\ \hline
\textbf{Usuario} & Inspector HSEQ \\ \hline
\textbf{Prioridad} & MUST \\ \hline
\textbf{Puntos Estimados de Esfuerzo} & 8 \\ \hline
\textbf{Descripción} & \parbox[t]{\linewidth}{\vspace{2pt}\begin{itemize}[leftmargin=1.5em, nosep, topsep=0pt]
\item \textbf{Como} Inspector HSEQ
\item \textbf{Quiero} habilitar una capa visual que muestre los permisos de trabajo activos (alturas, caliente, izaje) sobre el mapa de la planta
\item \textbf{Para} identificar rápidamente cruces de actividades peligrosas y evitar accidentes por interferencia entre contratistas.
\end{itemize}\vspace{2pt}} \\ \hline
\textbf{Observaciones} & Esta funcionalidad es crítica para la seguridad industrial (SIMOPS) y el cumplimiento de la norma ISO 45001. Más allá de la representación gráfica, el sistema implementa una validación de integridad (Detección de Objetos Huérfanos) que garantiza la correspondencia exacta entre los permisos digitales y los activos físicos en el modelo 3D, alertando sobre cualquier inconsistencia que pudiera comprometer la seguridad del personal. \\ \hline
\end{longtable}\normalsize
\noindent\textit{Criterios de Aceptación:}
\footnotesize\begin{longtable}{|p{0.28\textwidth}|p{0.32\textwidth}|p{0.34\textwidth}|}\hline
\textbf{Contexto} & \textbf{Acción} & \textbf{Resultado} \\ \hline
\parbox[t]{\linewidth}{Dado que el Inspector está visualizando el mapa general de planta (Nivel 1).} & \parbox[t]{\linewidth}{Cuando activa la capa ``Permisos de Trabajo''.} & \parbox[t]{\linewidth}{Entonces el sistema superpone iconos sobre los activos afectados, utilizando codificación redundante (Color + Forma): Trabajo en Caliente: Cuadrado Rojo. Trabajo en Alturas: Triángulo Azul. Espacio Confinado: Círculo Amarillo.} \\ \hline
\parbox[t]{\linewidth}{Dado que se muestra un icono de ``Trabajo en Caliente'' sobre la bomba P-202.} & \parbox[t]{\linewidth}{Cuando el Inspector hace clic sobre el icono.} & \parbox[t]{\linewidth}{Entonces se despliega una tarjeta compacta (sin tapar el contexto) con: Contratista, Hora Inicio/Fin, y Responsable HSEQ.} \\ \hline
\parbox[t]{\linewidth}{Dado que existe un Permiso de Trabajo activo asociado al Tag ``V-101''.} & \parbox[t]{\linewidth}{Cuando el sistema intenta renderizar el icono en el modelo 3D y detecta que el objeto ``V-101'' no existe en la geometría del modelo actual (Enlace Roto).} & \parbox[t]{\linewidth}{Entonces el sistema genera una alerta de ``Integridad de Seguridad Comprometida'' en el panel de notificaciones y muestra el permiso en una lista de ``Fallos de Visualización'' lateral, para que no pase desapercibido.} \\ \hline
\end{longtable}\normalsize
\vspace{1em}
\noindent\textbf{Ficha Técnica: VIS-009 - Navegación Jerárquica y Zoom Semántico}
\footnotesize\setlength{\tabcolsep}{4pt}\renewcommand{\arraystretch}{1.2}
\rowcolors{2}{tableBlue}{white}\begin{longtable}{|p{0.21\textwidth}|p{0.74\textwidth}|}\hline
\textbf{Número} & VIS-009 \\ \hline
\textbf{Usuario} & Supervisor de Mantenimiento \\ \hline
\textbf{Prioridad} & SHOULD \\ \hline
\textbf{Puntos Estimados de Esfuerzo} & 5 \\ \hline
\textbf{Descripción} & \parbox[t]{\linewidth}{\vspace{2pt}\begin{itemize}[leftmargin=1.5em, nosep, topsep=0pt]
\item \textbf{Como} Supervisor de Mantenimiento
\item \textbf{Quiero} hacer zoom en una zona específica de la planta usando el scroll del mouse
\item \textbf{Para} inspeccionar con detalle la ubicación de los activos en áreas densas
\end{itemize}\vspace{2pt}} \\ \hline
\textbf{Observaciones} & Esta funcionalidad de navegación avanzada mejora significativamente la experiencia de usuario (UX) mediante la implementación de zoom semántico. Para garantizar un rendimiento óptimo en dispositivos móviles y estaciones de trabajo, el sistema utiliza lógica de renderizado condicional (LOD), ajustando dinámicamente la complejidad de la geometría y la densidad de información según el nivel de acercamiento del usuario. \\ \hline
\end{longtable}\normalsize
\noindent\textit{Criterios de Aceptación:}
\footnotesize\begin{longtable}{|p{0.28\textwidth}|p{0.32\textwidth}|p{0.34\textwidth}|}\hline
\textbf{Contexto} & \textbf{Acción} & \textbf{Resultado} \\ \hline
\parbox[t]{\linewidth}{Dado que el Supervisor está viendo el Mapa General (Nivel 1).} & \parbox[t]{\linewidth}{Cuando hace zoom sobre el área de ``Calderas''.} & \parbox[t]{\linewidth}{Entonces el mapa cambia de mostrar un bloque genérico a mostrar los iconos individuales de los activos (Nivel 2) y sus etiquetas de estado.} \\ \hline
\parbox[t]{\linewidth}{Dado que el Supervisor ha alcanzado el nivel máximo de detalle.} & \parbox[t]{\linewidth}{Cuando intenta hacer más zoom.} & \parbox[t]{\linewidth}{Entonces la vista se bloquea suavemente (efecto elástico) para evitar que la cámara atraviese el suelo o entre en geometrías no renderizadas.} \\ \hline
\parbox[t]{\linewidth}{Dado que el Supervisor está inspeccionando una bomba específica.} & \parbox[t]{\linewidth}{Cuando presiona el botón ``Vista General'' o hace zoom out rápido.} & \parbox[t]{\linewidth}{Entonces la cámara se aleja hasta mostrar toda la unidad de proceso, restaurando la visión periférica de alarmas vecinas.} \\ \hline
\end{longtable}\normalsize
\vspace{1em}
\noindent\textbf{Ficha Técnica: VIS-010 - Visualización de Despiece (Vista Explosionada)}
\footnotesize\setlength{\tabcolsep}{4pt}\renewcommand{\arraystretch}{1.2}
\rowcolors{2}{tableBlue}{white}\begin{longtable}{|p{0.21\textwidth}|p{0.74\textwidth}|}\hline
\textbf{Número} & VIS-010 \\ \hline
\textbf{Usuario} & Técnico de Mantenimiento \\ \hline
\textbf{Prioridad} & COULD \\ \hline
\textbf{Puntos Estimados de Esfuerzo} & 8 \\ \hline
\textbf{Descripción} & \parbox[t]{\linewidth}{\vspace{2pt}\begin{itemize}[leftmargin=1.5em, nosep, topsep=0pt]
\item \textbf{Como} Técnico de Mantenimiento
\item \textbf{Quiero} visualizar la vista explosionada (despiece) del activo seleccionado en 3D
\item \textbf{Para} identificar la ubicación exacta de un componente interno antes de iniciar el desarme físico.
\end{itemize}\vspace{2pt}} \\ \hline
\textbf{Observaciones} & La vista explosionada (despiece) representa un nivel avanzado de visualización para facilitar intervenciones complejas. Debido a su alta demanda de recursos gráficos y modelado especializado, se clasifica como una funcionalidad opcional que requiere modelos 3D jerárquicos y optimizados para visualización web intensiva por hardware, ofreciendo una alternativa interactiva a los manuales técnicos tradicionales. \\ \hline
\end{longtable}\normalsize
\noindent\textit{Criterios de Aceptación:}
\footnotesize\begin{longtable}{|p{0.28\textwidth}|p{0.32\textwidth}|p{0.34\textwidth}|}\hline
\textbf{Contexto} & \textbf{Acción} & \textbf{Resultado} \\ \hline
\parbox[t]{\linewidth}{Dado que el Técnico ha seleccionado una bomba en el visor 3D.} & \parbox[t]{\linewidth}{Cuando selecciona la opción ``Ver Despiece''.} & \parbox[t]{\linewidth}{Entonces las partes del modelo se separan visualmente (animación de explosión) permitiendo ver los componentes internos como el eje y los rodamientos.} \\ \hline
\parbox[t]{\linewidth}{Dado que el activo está en modo ``Vista Explosionada''.} & \parbox[t]{\linewidth}{Cuando el Técnico selecciona un rodamiento específico.} & \parbox[t]{\linewidth}{Entonces el sistema resalta el componente y muestra una etiqueta con su nombre y código de parte vinculado al Inventario.} \\ \hline
\end{longtable}\normalsize
\vspace{1em}
\noindent\textbf{Ficha Técnica: VIS-011 - Visualización de Rutas de Aislamiento (LOTO) y Bloqueo Activo}
\footnotesize\setlength{\tabcolsep}{4pt}\renewcommand{\arraystretch}{1.2}
\rowcolors{2}{tableBlue}{white}\begin{longtable}{|p{0.21\textwidth}|p{0.74\textwidth}|}\hline
\textbf{Número} & VIS-011 \\ \hline
\textbf{Usuario} & Técnico de Mantenimiento \\ \hline
\textbf{Prioridad} & MUST \\ \hline
\textbf{Puntos Estimados de Esfuerzo} & 13 \\ \hline
\textbf{Descripción} & \parbox[t]{\linewidth}{\vspace{2pt}\begin{itemize}[leftmargin=1.5em, nosep, topsep=0pt]
\item \textbf{Como} Técnico de Mantenimiento
\item \textbf{Quiero} visualizar la ruta de aislamiento de energía (LOTO) sobre el modelo 3D y que el sistema bloquee el inicio del trabajo si detecta energía activa.
\item \textbf{Para} asegurar que la intervención se realice bajo condiciones de Energía Cero, previniendo accidentes laborales.
\end{itemize}\vspace{2pt}} \\ \hline
\textbf{Observaciones} & Esta funcionalidad constituye el diferenciador tecnológico clave del Gemelo Digital, integrando visualización 3D avanzada con lógica de seguridad crítica en tiempo real (Energía Cero). El sistema combina la trazabilidad gráfica de puntos de aislamiento (Lockout/Tagout) con un mecanismo de bloqueo activo por software que consume telemetría IoT, garantizando que el inicio de cualquier intervención cumpla rigurosamente con los protocolos de seguridad industrial más exigentes. \\ \hline
\end{longtable}\normalsize
\noindent\textit{Criterios de Aceptación:}
\footnotesize\begin{longtable}{|p{0.28\textwidth}|p{0.32\textwidth}|p{0.34\textwidth}|}\hline
\textbf{Contexto} & \textbf{Acción} & \textbf{Resultado} \\ \hline
\parbox[t]{\linewidth}{Dado que el Técnico tiene seleccionada una Orden de Trabajo en su tablet.} & \parbox[t]{\linewidth}{Cuando activa la función ``Ver Ruta LOTO''.} & \parbox[t]{\linewidth}{Entonces el Gemelo Digital resalta en el modelo 3D el activo, las fuentes de energía y los puntos de bloqueo físicos (ej. válvulas, interruptores) conectados mediante líneas de trazabilidad.} \\ \hline
\parbox[t]{\linewidth}{Dado que se muestran las líneas de conexión LOTO.} & \parbox[t]{\linewidth}{Cuando el Técnico selecciona el punto de corte sugerido (ej. Breaker).} & \parbox[t]{\linewidth}{Entonces el sistema debe mostrar la etiqueta del tablero (Tag) que coincide con la realidad física.} \\ \hline
\parbox[t]{\linewidth}{Dado que el sistema detecta energía activa.} & \parbox[t]{\linewidth}{Cuando el técnico intenta presionar ``Iniciar Trabajo''.} & \parbox[t]{\linewidth}{Entonces el sistema bloquea la acción y muestra una alerta de peligro.} \\ \hline
\parbox[t]{\linewidth}{Dado que los sensores confirman que no hay energía activa.} & \parbox[t]{\linewidth}{Cuando el técnico confirma el bloqueo físico.} & \parbox[t]{\linewidth}{Entonces el sistema permite cambiar el estado de la WO a ``En Ejecución''.} \\ \hline
\parbox[t]{\linewidth}{Dado que el sistema no puede establecer comunicación con los sensores IoT de la zona.} & \parbox[t]{\linewidth}{Cuando el técnico intenta iniciar el trabajo o confirmar el aislamiento.} & \parbox[t]{\linewidth}{Entonces el sistema bloquea preventivamente la operación, informando que no se puede garantizar la Energía Cero por falta de datos en tiempo real.} \\ \hline
\end{longtable}\normalsize
\vspace{1em}
\noindent\textbf{Ficha Técnica: VIS-012 - Configuración de Jerarquía de Componentes 3D}
\footnotesize\setlength{\tabcolsep}{4pt}\renewcommand{\arraystretch}{1.2}
\rowcolors{2}{tableBlue}{white}\begin{longtable}{|p{0.21\textwidth}|p{0.74\textwidth}|}\hline
\textbf{Número} & VIS-012 \\ \hline
\textbf{Usuario} & Ingeniero de Proyectos \\ \hline
\textbf{Prioridad} & COULD \\ \hline
\textbf{Puntos Estimados de Esfuerzo} & 5 \\ \hline
\textbf{Descripción} & \parbox[t]{\linewidth}{\vspace{2pt}\begin{itemize}[leftmargin=1.5em, nosep, topsep=0pt]
\item \textbf{Como} Ingeniero de Proyectos
\item \textbf{Quiero} definir la jerarquía de componentes de un activo (ej. Bomba \$\textbackslash\{\}to\$ Motor \$\textbackslash\{\}to\$ Rodamiento) y asociar un modelo 3D a cada nivel
\item \textbf{Para} construir la vista explosionada del Gemelo Digital y permitir que los técnicos accedan a la información del componente específico en campo.
\end{itemize}\vspace{2pt}} \\ \hline
\textbf{Observaciones} & Esta herramienta administrativa permite la configuración granular de la jerarquía de activos, facilitando la construcción del Gemelo Digital multinivel. Arquitectónicamente, se sugiere seguir el patrón Sidecar para el almacenamiento pesado (modelos 3D en almacenamiento de blob/objetos y metadatos en backend central), garantizando alta escalabilidad. Aunque el MVP operará con una estructura simplificada, esta funcionalidad asegura la viabilidad a futuro. \\ \hline
\end{longtable}\normalsize
\noindent\textit{Criterios de Aceptación:}
\footnotesize\begin{longtable}{|p{0.28\textwidth}|p{0.32\textwidth}|p{0.34\textwidth}|}\hline
\textbf{Contexto} & \textbf{Acción} & \textbf{Resultado} \\ \hline
\parbox[t]{\linewidth}{Dado que el Ingeniero está en la página de un activo ``Bomba P-101'' que no tiene componentes definidos.} & \parbox[t]{\linewidth}{Cuando hace clic en ``Añadir Componente'', le da el nombre ``Motor Eléctrico'' y el código ``ME-01''.} & \parbox[t]{\linewidth}{Entonces ``Motor Eléctrico'' debe aparecer como un sub-ítem anidado debajo de la ``Bomba P-101'' en el árbol de activos.} \\ \hline
\parbox[t]{\linewidth}{Dado que el Ingeniero está editando las propiedades del componente ``Motor Eléctrico''.} & \parbox[t]{\linewidth}{Cuando utiliza la opción ``Subir Modelo 3D'' y adjunta un archivo motor.step.} & \parbox[t]{\linewidth}{Entonces el sistema debe asociar ese archivo 3D al componente y mostrar una vista previa en miniatura.} \\ \hline
\parbox[t]{\linewidth}{Dado que el Técnico está viendo el activo ``Bomba P-101'' en la app móvil.} & \parbox[t]{\linewidth}{Cuando hace clic en la vista explosionada y selecciona el ``Motor Eléctrico''.} & \parbox[t]{\linewidth}{Entonces la app debe mostrar la información y el modelo 3D específico del ``Motor Eléctrico'', no el de la bomba completa.} \\ \hline
\parbox[t]{\linewidth}{Dado que el activo ``Bomba P-101'' ya tiene un componente con el código ``ME-01''.} & \parbox[t]{\linewidth}{Cuando el Ingeniero intenta crear otro componente con el mismo código ``ME-01''.} & \parbox[t]{\linewidth}{Entonces el sistema debe mostrar un mensaje de error ``El código de componente ya existe en este activo'' y no debe crear el duplicado.} \\ \hline
\parbox[t]{\linewidth}{Dado que el componente ``Rodamiento R-01'' fue creado por error dentro del ``Motor A''.} & \parbox[t]{\linewidth}{Cuando el Ingeniero arrastra el componente hacia el ``Motor B'' en el árbol de jerarquía.} & \parbox[t]{\linewidth}{Entonces el sistema debe actualizar la relación padre-hijo, moviendo el rodamiento y todo su historial al nuevo padre.} \\ \hline
\end{longtable}\normalsize
\vspace{1em}